\chapter{Resultados Detalhados do Treinamento Exploratório}\label{ap:resultados_exploratorios_detalhados_rev} % Label revisado

Este apêndice apresenta os resultados detalhados para cada configuração de modelo e hiperparâmetros avaliada durante a fase de treinamento exploratório. Os valores correspondem à média do desempenho no conjunto de teste ao longo dos 5 \textit{folds} de validação cruzada (conforme descrito na \Secao{sec:res_config_exp_exploratoria_subamostra}). Os pesos da função de perda combinada ($\alpha, \beta, \gamma$) foram fixados em 0.5, 0.25 e 0.25, respectivamente, para todas as execuções.

\begin{landscape} % Inicia o ambiente de paisagem
\begin{table}[!htp]
\centering
\caption[Resultados detalhados do treinamento exploratório por modelo e configuração.]{Métricas de avaliação e hiperparâmetros para os diferentes modelos testados na fase exploratória. LR: Taxa de Aprendizado; Freeze Epo: Número de épocas com \textit{encoder} congelado (0 indica que não houve congelamento); Weight D: Decaimento de Peso. Células destacadas indicam valores notáveis ou melhores para a respectiva métrica na configuração específica.}
\label{tab:ap_resultados_exploratorios_completos_rev} % Label revisado para a tabela
\resizebox{\linewidth}{!}{% Ajusta a tabela para a largura da página em modo paisagem
\sisetup{output-decimal-marker={,}} % Configura o siunitx para usar vírgula como separador decimal
\begin{tabular}{l S[table-format=1.4] S[table-format=1.4] S[table-format=1.4] S[table-format=1.4] S[table-format=1.4] S[table-format=1.4] c S[table-format=1.6] c S[table-format=1.6] l l c S[table-format=3.1,table-parse-only,table-figures-decimal=1]}
% Ajuste os table-format conforme o número de dígitos inteiros e decimais de cada coluna
\toprule
\textbf{Test Loss} & \textbf{Accuracy} & \textbf{Dice} & \textbf{IoU} & \textbf{TPR} & \textbf{TNR} & \textbf{Precision} & \textbf{Dropout} & \textbf{LR} & \textbf{Freeze Epo} & \textbf{Weight D} & \textbf{MODEL} & \textbf{ENCODER} & \textbf{Frozen?} & \textbf{PARAMETERS (M)} \\ 
\midrule
0,4383 & 0,89488 & \cellcolor{highlightcolor}0,78984 & \cellcolor{highlightcolor}0,65454 & \cellcolor{highlightcolor}0,7484  & 0,94564 & 0,84068 & 0,2 & 0,00001 & 0 & 0,00001 & SWIN & large & NO & 197 \\
0,4389 & \cellcolor{highlightcolor}0,89654 & 0,78898 & 0,65392 & 0,74564 & 0,94598 & 0,84696 & 0,2 & 0,00001 & 0 & 0,00001 & SWIN-LARGE-DICE & large & NO & 197 \\
0,4403 & 0,89402 & 0,78036 & 0,6418  & 0,73384 & 0,94838 & 0,84006 & 0,2 & 0,00001 & 0 & 0,00001 & SWIN & base & NO & 88 \\
\cellcolor{highlightcolor} 
0,4307 & 0,88688 & 0,76954 & 0,62656 & 0,72518 & 0,94154 & 0,83382 & 0,2 & 0,00001 & 0 & 0,00001 & DeepLabV3Plus & resnet101 & NO & 42 \\
0,43866 & 0,88504 & 0,76356 & 0,61916 & 0,72738 & 0,93752 & 0,81838 & 0,2 & 0,00001 & 0 & 0,00001 & DeepLabV3Plus & resnet152 & NO & 58 \\
0,46184 & 0,88078 & 0,75426 & 0,60754 & 0,71322 & 0,93518 & 0,80144 & 0,2 & 0,00001 & 0 & 0,00001 & UNET & inceptionresnetv2 & NO & 54 \\
0,47734 & 0,8764  & 0,75416 & 0,6072  & 0,74456 & 0,91614 & 0,78756 & 0,2 & 0,00001 & 0 & 0,00001 & DPT & tu-vit\_base & NO & 102,6 \\
0,47438 & 0,87478 & 0,74696 & 0,59752 & 0,70006 & 0,9362  & 0,8043  & 0,2 & 0,00001 & 0 & 0,00001 & UNET++ & resnet152 & NO & 58 \\
0,47414 & 0,87752 & 0,73716 & 0,5857  & 0,65482 & \cellcolor{highlightcolor}0,95358 & \cellcolor{highlightcolor}0,8481  & 0,2 & 0,00001 & 0 & 0,00001 & UNET++ & resnet101 & NO & 42 \\
0,50214 & 0,87152 & 0,7319  & 0,5789  & 0,68738 & 0,93312 & 0,80086 & 0,2 & 0,00001 & 0 & 0,00001 & FPN & resnet152 & NO & 58 \\
0,47486 & 0,8717  & 0,7318  & 0,58066 & 0,7095  & 0,91814 & 0,78104 & 0,2 & 0,00001 & 0 & 0,00001 & Segformer & mit\_b5 & NO & 81 \\
0,48848 & 0,86906 & 0,72528 & 0,57176 & 0,65276 & 0,94772 & 0,8197  & 0,2 & 0,00001 & 0 & 0,00001 & Manet & resnet152 & NO & 58 \\
0,51508 & 0,86408 & 0,7212  & 0,56654 & 0,6864  & 0,91782 & 0,77898 & 0,2 & 0,00001 & 0 & 0,00001 & FPN & resnet101 & NO & 42 \\
0,48046 & 0,85836 & 0,69832 & 0,5408  & 0,6665  & 0,91124 & 0,77158 & 0,2 & 0,00001 & 0 & 0,00001 & Segformer & mit\_b3 & NO & 44 \\
0,47052 & 0,83882 & 0,6759  & 0,5135  & 0,6675  & 0,894   & 0,70202 & 0   & 0,0001  & 2 & 0,0001  & SWIN & base & YES & 88 \\
0,50084 & 0,8374  & 0,67392 & 0,5093  & 0,65018 & 0,90064 & 0,7044  & 0,2 & 0,0001  & 2 & 0,0001  & DeepLabV3Plus & resnet50 & YES & 23 \\
0,49988 & 0,83448 & 0,66566 & 0,49982 & 0,63482 & 0,90586 & 0,70992 & 0,2 & 0,0001  & 2 & 0,0001  & DeepLabV3Plus & resnet101 & YES & 42 \\
0,47854 & 0,8302  & 0,65154 & 0,48806 & 0,64228 & 0,88522 & 0,68114 & 0   & 0,0001  & 2 & 0,0001  & SWIN & tiny & YES & 28 \\
0,47874 & 0,8338  & 0,65146 & 0,48552 & 0,61064 & 0,90558 & 0,71626 & 0,2 & 0,0001  & 2 & 0,0001  & UNET++ & resnet50 & YES & 23 \\
0,52454 & 0,82736 & 0,64458 & 0,47776 & 0,60494 & 0,90368 & 0,69514 & 0,2 & 0,0001  & 2 & 0,0001  & FPN & resnet50 & YES & 23 \\
0,49696 & 0,81952 & 0,62826 & 0,46092 & 0,59736 & 0,89222 & 0,67232 & 0,2 & 0,0001  & 2 & 0,0001  & UNET & resnet50 & YES & 23 \\
0,51134 & 0,80502 & 0,60634 & 0,43924 & 0,60028 & 0,86838 & 0,62504 & 0,2 & 0,0001  & 2 & 0,0001  & Segformer & mit\_b2 & YES & 24 \\
\bottomrule
\end{tabular}
}
\end{table}
\end{landscape} % Termina o ambiente de paisagem

\chapter{Esquema do Banco de Dados de Gerenciamento}\label{ap:esquema_banco_dados}

Este apêndice detalha a estrutura da tabela principal do banco de dados \textit{SQLite}, cuja nomenclatura segue o padrão dinâmico \texttt{DATABASE\_\{TAG\}} (onde \texttt{\{TAG\}} representa o identificador do experimento, \eg, \texttt{DATABASE\_V1}). Esta tabela é responsável por orquestrar a fila de processamento, armazenar os caminhos dos arquivos e registrar os metadados de execução. A \Tabela{tab:ap_esquema_database} apresenta as colunas, seus tipos de dados conforme implementação em \textit{SQLite} e a descrição de sua finalidade.

\begin{table}[!htb]
\centering
\caption[Esquema da tabela de gerenciamento DATABASE\_\{TAG\}.]{Estrutura detalhada da tabela \texttt{DATABASE\_\{TAG\}} do banco de dados de gerenciamento.}
\label{tab:ap_esquema_database}
% Ajuste a largura das colunas p{...} conforme necessário
\begin{tabular}{l l p{6.5cm}} 
\toprule
\textbf{Coluna} & \textbf{Tipo} & \textbf{Descrição} \\
\midrule
ID & \texttt{INTEGER (PK)} & Identificador único sequencial para cada registro (Chave Primária). \\
STATUS & \texttt{TEXT} & Estado atual do processamento (\eg, \texttt{'TO BE PROCESSED'}, \texttt{'COMPLETED'}, \texttt{'FAILED'}). Valor padrão: \texttt{'TO BE PROCESSED'}. \\
IMAGEPATH & \texttt{TEXT} & Caminho absoluto para o arquivo de imagem digital (\eg, \texttt{.svs}, \texttt{.tif}). \\
ANNOTATIONPATH & \texttt{TEXT} & Caminho absoluto para o arquivo de anotação correspondente (\eg, \texttt{.xml}). \\
LastUpdate & \texttt{TIMESTAMP} & Carimbo de data/hora da última modificação do registro. \\
CANCER\_QTD & \texttt{INTEGER} & Contagem total de \textit{patches} da classe "Câncer" gerados. \\
NON\_CANCER\_QTD & \texttt{INTEGER} & Contagem total de \textit{patches} da classe "Não Câncer" gerados. \\
VALID\_IMAGE & \texttt{INTEGER} & Flag binário indicando se a imagem passou na validação de integridade inicial. \\
CANCER\_COLOR & \texttt{TEXT} & Código de cor (nome ou hexadecimal) utilizado nas anotações para identificar a classe positiva. \\
NOT\_CANCER\_COLOR & \texttt{TEXT} & Código de cor (nome ou hexadecimal) utilizado nas anotações para identificar a classe negativa. \\
PROCESSINGTIME\_MINUTES & \texttt{REAL} & Tempo total de execução do algoritmo de \textit{patching} para a lâmina (em minutos). \\
PATIENT & \texttt{TEXT} & Identificador único do paciente (Restrição \texttt{UNIQUE}). \\
COMMENTS & \texttt{TEXT} & Logs de execução, mensagens de erro ou observações geradas pelo \textit{script}. \\
WINDOW\_SIZE & \texttt{INTEGER} & Tamanho da janela de recorte do \textit{patch} (em pixels). \\
STRIDE & \texttt{INTEGER} & Passo de deslocamento da janela (\textit{stride}) durante o recorte (em pixels). \\
MATCH\_PERCENTAGE & \texttt{TEXT} & Parâmetro de limiar de sobreposição entre a anotação e o \textit{patch} para classificação. \\
TISSUE\_PERCENTAGE & \texttt{TEXT} & Parâmetro de limiar mínimo de área de tecido para considerar o \textit{patch} válido. \\
\bottomrule
\end{tabular}
\end{table}